%=========================== RESUMO ===========================
\begingroup
\begin{center}{\bfseries RESUMO}\end{center}
\vspace{0.3cm}

\noindent No Brasil, estima-se que mais de 10 milhões de pessoas possuam algum
grau de deficiência auditiva, e a comunicação entre essa população e pessoas
ouvintes depende, na maior parte dos contextos cotidianos, da presença de um
intérprete de Língua Brasileira de Sinais (Libras) — recurso escasso frente à
demanda real. Este trabalho apresenta o LibrasAI, um sistema de reconhecimento
automático de gestos dinâmicos de Libras em tempo real, desenvolvido a partir de
técnicas de visão computacional e aprendizado profundo. O sistema captura vídeo
de uma webcam convencional, extrai 21 \textit{landmarks} de cada mão por meio da
biblioteca MediaPipe, normaliza espacialmente essas coordenadas e as processa
através de uma rede neural híbrida CNN + LSTM, capaz de classificar sequências
temporais de 50 \textit{frames} e convertê-las em texto, formando palavras e
frases de maneira contínua. A pesquisa se classifica como pesquisa aplicada do
problema, seguida do desenvolvimento de uma aplicação funcional. O protocolo de
coleta de dados envolveu 10 participantes, com no mínimo 10 repetições por gesto,
em ambientes de fundo neutro e iluminação controlada. Resultados preliminares,
obtidos com um \textit{dataset} inicial de 2 classes de gestos, indicam acurácia
de 97,4\%, F1-score macro de 0,97 e coeficiente Kappa de Cohen de 0,94 — valores
condizentes com a literatura da área, ainda que obtidos sobre uma base de dados
reduzida. São discutidas as limitações do trabalho, sobretudo quanto ao tamanho
do \textit{dataset} e à ausência de validação com usuários externos ao
treinamento, e propostas as etapas seguintes de expansão do vocabulário
reconhecido e de validação com a comunidade surda.

\vspace{0.4cm}
\noindent\textbf{Palavras-chave:} Língua Brasileira de Sinais; Libras;
Reconhecimento de gestos; Visão computacional; Aprendizado profundo; CNN; LSTM;
MediaPipe.

\vspace{1.0cm}
\begin{center}{\bfseries ABSTRACT}\end{center}
\vspace{0.3cm}

\noindent In Brazil, more than 10 million people are estimated to have some
degree of hearing impairment, and communication between this population and
hearing people depends, in most everyday contexts, on the presence of a Brazilian
Sign Language (Libras) interpreter — a scarce resource relative to actual demand.
This work presents LibrasAI, a real-time automatic recognition system for dynamic
Libras gestures, built upon computer vision and deep learning techniques. The
system captures video from a conventional webcam, extracts 21 landmarks per hand
using the MediaPipe library, spatially normalizes these coordinates, and
processes them through a hybrid CNN + LSTM neural network capable of classifying
temporal sequences of 50 frames and converting them into text, continuously
forming words and sentences. The research is classified as applied
problem-oriented research followed by the development of a functional
application. The data collection protocol involved 10 participants, with at least
10 repetitions per gesture, in neutral-background, controlled-lighting
environments. Preliminary results, obtained with an initial dataset of 2 gesture
classes, indicate 97.4\% accuracy, a macro F1-score of 0.97, and a Cohen's Kappa
of 0.94 — values consistent with the literature, albeit obtained over a small
dataset. The work's limitations are discussed, particularly regarding dataset
size and the lack of validation with users external to training, along with the
next steps for expanding the recognized vocabulary and validating the system with
the deaf community.

\vspace{0.4cm}
\noindent\textbf{Keywords:} Brazilian Sign Language; Libras; Gesture recognition;
Computer vision; Deep learning; CNN; LSTM; MediaPipe.
\endgroup
\newpage
